\subsection{Aproximação longa}
\label{sec:aproximacao_longa}

A fase de aproximação longa inicia-se após o estabelecimento da órbita de fase, quando o perseguidor se encontra a uma distância da ordem de centenas de metros em relação ao alvo no referencial \gls{lvlh}.
Nessa condição, a separação relativa é suficientemente grande para que pequenas imprecisões de manobra não representem risco imediato de colisão, mas já pequena o bastante para que a dinâmica relativa seja realizado entre as duas espaçonaves.

O critério de entrada na aproximação longa é superior às dimensões da região de captura a ser utilizada na aproximação curta.
Além da posição, são especificados limites máximos de velocidade relativa em cada componente, de forma a garantir que o perseguidor chegue à vizinhança do alvo com velocidades controladas e compatíveis com a capacidade de correção nas fases subsequentes.

Durante a aproximação longa, o perseguidor reduz gradualmente a distância relativa para evitar que se aproxime do alvo com velocidades elevadas, bem como regiões em que a perda de controle possa levar a colisão ou afastamento não planejado, o que comprometeria a missão.

Quando a posição relativa do perseguidor se encontra dentro de uma região e as velocidades relativas em cada eixo do \gls{lvlh} estão abaixo de valores determinados, considera-se que as condições para a aproximação longa foram atendidas.
O estado relativo atingido ao final desta fase serve como condição inicial nominal para a aproximação curta.
